\documentclass[11pt]{article}

\usepackage[margin=2.5cm]{geometry}
\usepackage{times}
\usepackage{amsmath}
\usepackage{booktabs}
\usepackage{graphicx}
\usepackage{pgfplots}
\pgfplotsset{compat=1.17}
\usepackage[numbers,sort&compress]{natbib}
\usepackage{hyperref}
\hypersetup{colorlinks=true,linkcolor=blue,citecolor=blue,urlcolor=blue}

% Numbers generated by the evaluation harness (scripts/evaluate.ts).
% Run `npm run evaluate` to (re)generate data/*.tex before building.
\newcommand{\numreplicates}{30}
\newcommand{\baseseed}{20260713}
\newcommand{\numfeatures}{10}
\newcommand{\trainepisodes}{6}
\newcommand{\healthyrecall}{98.7}
\newcommand{\illrecall}{99.9}
\newcommand{\worstfp}{0.01}
\newcommand{\socialworstfp}{0.01}
\newcommand{\illlatsocial}{2.4}
\newcommand{\illlattelemetry}{2.4}
\newcommand{\illlatgain}{0.0}


\title{\vspace{-1em}Herd-Relative and Social-Network Anomaly Detection for\\
Livestock Wearables: An In-Silico Study}

\author{HerdLink Project\\\small A reproducible simulation benchmark and detection method}
\date{\today}

\begin{document}
\maketitle

\begin{abstract}
\noindent
Health-monitoring features on GPS livestock collars almost universally raise
alerts by comparing each animal's activity to a \emph{fixed per-animal
baseline}. This design is fragile to herd-wide covariates---weather, feed, time
of day---which shift every animal at once and either trigger false alarms or
force detuning that misses real cases. We present a detection method that (i)
scores every animal \emph{relative to its herd and to its own recent history},
and (ii) mines the herd's proximity-derived \emph{social network} for early
signs of illness, chiefly social withdrawal. Because no public, labelled dataset
of collar telemetry with ground-truth health states exists, we release a
controllable generative herd simulator and evaluate on it. \textbf{This is an
in-silico proof of concept: no real-herd validation is performed, and the
results below are evidence about the method within the simulation, not field
performance.} Across \numreplicates{} seeded replicates, a cost-sensitive
multinomial logistic classifier over \numfeatures{} herd-, self-, and
social-relative features attains \healthyrecall\%{} held-out recall on healthy
animals while raising at most \socialworstfp{} false alerts per cow-day under
forced heatwave, snow, and rain regimes---the central evidence for herd-relative
normalisation. A pre-specified feature-family ablation returns a deliberately
reported \emph{null} result: self- and social-relative features do not shorten
detection latency in this benchmark, because illness is already identified
rapidly from fever and rumination and spatial withdrawal is already captured by a
herd-centroid-distance feature. The social features add interpretable network
structure without inflating false alarms, and we identify the conditions
(noisier field telemetry) under which they would be expected to contribute. We
release the simulator, detector, and evaluation harness to support replication
and, ultimately, field validation.
\end{abstract}

\section{Introduction}
Virtual-fencing collars place a continuously-sampling sensor on every animal in
a herd, and health alerting is among their most valuable features
\citep{berckmans2017plf,rutten2013economic}. The dominant alerting approach
monitors an individual's activity and rumination against its own fixed baseline
and flags deviations \citep{stangaferro2016ketosis}. This has a structural
weakness: behaviour is driven by covariates that move the \emph{entire herd}
simultaneously. A heatwave raises every body temperature; a storm halves every
grazing bout. Fixed thresholds must then choose between firing on these
herd-wide shifts (alert fatigue, which erodes trust and ultimately use) or being
detuned until genuine cases are missed.

We argue two things follow from taking the herd, rather than the animal, as the
unit of normalisation. First, scoring each animal \emph{relative to the current
herd distribution} makes detection invariant to herd-wide covariates: when
everyone shifts together the distribution moves with them and no individual
stands out. Second, a collar network is implicitly a \emph{social sensor}: every
position report is a proximity observation, and proximity over time reveals herd
social structure. Social withdrawal is a well-documented early sign of illness in
cattle \citep{weary2009painful,proudfoot2012social}, and it is latent in data
collars already collect.

\paragraph{Contributions.}
(1) A detection method combining herd-relative, self-relative, and
social-network features under cost-sensitive calibration. (2) A bias-corrected
online estimator of proximity tie strength suitable for streaming collar data.
(3) A controllable, fully labelled generative herd benchmark and a reproducible
evaluation harness. (4) A pre-specified ablation isolating the marginal
value---and, in this benchmark, the current redundancy---of self-relative and
social features, reported as a null result rather than suppressed.

\paragraph{Scope and honesty statement.}
There is, to our knowledge, no public dataset pairing collar telemetry with
veterinary ground truth at the scale this study needs. We therefore evaluate
entirely in simulation. The simulator encodes behavioural assumptions that the
detector then exploits; the results are best read as a demonstration that the
method \emph{can} recover injected conditions robustly, and as a hypothesis
generator for field trials---not as evidence of real-world accuracy. Section
\ref{sec:limitations} treats this at length.

\section{Related work}
\textbf{Activity and rumination monitoring.} Commercial and research systems
detect ketosis, mastitis, lameness, and oestrus from accelerometer-derived
activity and rumination time \citep{stangaferro2016ketosis,rutten2013economic}.
These typically use per-animal baselines and thresholds. \textbf{Behavioural
sickness signs.} Reduced feeding, altered activity, and social withdrawal are
established behavioural indicators of ill health
\citep{weary2009painful,proudfoot2012social}. \textbf{Animal social network
analysis.} Proximity loggers and observation have been used to reconstruct
livestock social networks and show they are structured and temporally stable
\citep{boyland2016social,croft2008exploring,farine2015constructing}. \textbf{
Anomaly detection with population baselines.} Scoring individuals against a
population rather than a fixed threshold is a standard idea in anomaly detection
\citep{chandola2009anomaly}; our contribution is to combine herd-relative,
self-relative, and network features for this domain and to calibrate for the
asymmetric cost of false alarms \citep{king2001logistic}.

\section{Generative herd model}
The simulator (a browser/Node TypeScript implementation) advances a herd of
grazing cattle in a bounded paddock. Each animal occupies one of four behavioural
states---grazing, walking, resting, ruminating---drawn from a time-of-day
distribution with dawn/dusk grazing peaks, blended with the herd's current state
mix to reproduce behavioural synchrony (allelomimicry). Weather (ambient
temperature, wind, rain/snow, cloud) evolves via rolling fronts and modulates the
state distribution and herd cohesion; a diurnal temperature cycle is
superimposed. Each animal holds a persistent spatial offset within the herd,
yielding stable proximity structure.

\paragraph{Telemetry.} Every five simulated minutes each collar emits a sample
$(\mathbf{p}, v, b, T, \rho)$: position, mean speed, behaviour class, body
temperature, and rumination fraction. A 24\,h ring buffer (288 samples) is
retained per animal. The detector observes only this stream.

\paragraph{Injected conditions.} Three ground-truth conditions perturb an
animal's pace, state distribution, body temperature, and spatial cohesion:
\emph{lameness} (markedly reduced pace, more lying, mild isolation);
\emph{illness} (fever $+1.4^\circ$C, collapsed rumination/grazing, strong
withdrawal); and \emph{oestrus} (elevated activity and walking, reduced rest).
These parameters define the labelled benchmark; they are not claims of
quantitative biological fidelity.

\section{Detection method}
\label{sec:method}

\subsection{Features}
For each animal at each 5-minute step we compute \numfeatures{} features in three
families, each standardised to a $z$-score against the current herd distribution
(with a floor on the standard deviation to avoid division by near-zero in a
perfectly synchronised herd):
\begin{itemize}\itemsep2pt
  \item \textbf{Herd-relative (5):} mean speed, walking fraction, rumination,
  body temperature, and distance from the herd centroid.
  \item \textbf{Self-relative (3):} the change in speed, walking fraction, and
  rumination between the last 2\,h and the animal's own 2--24\,h baseline. These
  separate a constitutionally slow animal from one that has \emph{become} slow.
  \item \textbf{Social (2):} network strength (weighted degree) and its change
  versus the animal's own baseline. These separate a constitutionally peripheral
  animal from one that is \emph{withdrawing}.
\end{itemize}

\subsection{Online social-tie estimation}
Pairs within a proximity radius (15\,m) accumulate association weight under an
exponential decay $\lambda$ (half-life $\approx 4$\,h): each 5-minute step, all
weights are multiplied by $\lambda$ and co-present pairs incremented by
$1-\lambda$, so a weight converges to the fraction of recent time spent together.
An exponential accumulator is biased low before it reaches steady state, so we
divide by the standard correction $1-\lambda^{t}$ at tick $t$
\citep{kingma2015adam}, giving an unbiased tie-strength estimate from the first
minutes of data. Weighted degree over the corrected ties gives each animal's
\emph{strength}.

\subsection{Classifier and alerting}
A multinomial logistic regression maps the feature vector to
$P(\text{healthy}/\text{ill}/\text{lame}/\text{oestrus})$. Training minimises
cross-entropy by full-batch gradient descent with $L_2$ regularisation and a
$3\times$ weight on healthy examples, reflecting the asymmetric on-farm cost of
false alarms \citep{king2001logistic}. At inference, an anomaly score
$s = 2.2\,(1-P(\text{healthy}))$ is smoothed by an EWMA ($\approx$1\,h) so alerts
require a \emph{sustained} signal; $s>2.0$ raises an alert and $s<1.0$ resolves
it. The predicted non-healthy class provides a human-readable suspected cause,
and the contributing $z$-scores provide a plain-language explanation.

\section{Experimental setup}
All randomness derives from a fixed base seed (\baseseed), making every number
reproducible via \texttt{npm run evaluate}. We run \numreplicates{} replicates.
Each replicate harvests \trainepisodes{} training episodes and 2 held-out test
episodes (60 animals, 36\,h, two animals per condition injected at 18\,h, samples
labelled positive only after a 2.5\,h onset gate). Ablation trains three models
per replicate---\emph{telemetry} (herd-relative only), \emph{self}
(+self-relative), \emph{social} (full \numfeatures{})---by masking feature
columns, so all share the live inference path. We report:

\begin{description}\itemsep2pt
  \item[A. Classification] held-out per-class precision/recall.
  \item[B. False positives] alerts per cow-day on \emph{healthy} herds run
  through the live detector under four weather regimes.
  \item[C. Detection latency] time from injection to first alert on the affected
  animal (14\,h warm-up, 12\,h observation window).
  \item[D. Ablation] B and C across the three feature families.
\end{description}
Means are reported with 95\% confidence intervals (normal approximation);
latencies as medians with interquartile range across replicates. On a laptop the
full harness completes in a few minutes.

\section{Results}
\label{sec:results}

\paragraph{Classification (A).} Held-out precision and recall
(Table~\ref{tab:cls}) show strong separation of illness and oestrus and weaker
recall on lameness, whose behavioural signature (reduced pace) overlaps most with
natural between-animal variation.

\begin{table}[t]\centering
\caption{Held-out classification, full model. Mean over \numreplicates{}
replicates (\%).}
\label{tab:cls}
\begin{tabular}{lrr}
\toprule
Class & Precision & Recall \\
\midrule
healthy & 93.3 & 98.7 \\
ill & 99.6 & 99.9 \\
lame & 94.8 & 77.8 \\
oestrus & 97.8 & 93.4 \\
\bottomrule
\end{tabular}

\end{table}

\paragraph{False positives under weather (B).} Table~\ref{tab:fp} reports alerts
per cow-day on \emph{healthy} herds. Rates remain near zero even under a
sustained forced heatwave and snowstorm---the central evidence that
herd-relative scoring absorbs herd-wide covariate shifts. Figure~\ref{fig:fp}
shows the same across ablation levels.

\begin{table}[t]\centering
\caption{False alerts per cow-day on healthy herds (mean\,$\pm$\,95\% CI), by
weather regime and feature family.}
\label{tab:fp}
\begin{tabular}{lrrr}
\toprule
Regime & Telemetry & +Self & +Social \\
\midrule
auto & 0.01\,$\pm$\,0.00 & 0.01\,$\pm$\,0.00 & 0.01\,$\pm$\,0.00 \\
heatwave & 0.01\,$\pm$\,0.00 & 0.01\,$\pm$\,0.00 & 0.01\,$\pm$\,0.00 \\
snow & 0.01\,$\pm$\,0.00 & 0.01\,$\pm$\,0.00 & 0.01\,$\pm$\,0.00 \\
rain & 0.01\,$\pm$\,0.00 & 0.01\,$\pm$\,0.01 & 0.01\,$\pm$\,0.00 \\
\bottomrule
\end{tabular}

\end{table}

\begin{figure}[t]\centering
\begin{tikzpicture}
\begin{axis}[
  width=0.7\linewidth, height=5cm,
  ybar, bar width=5pt,
  ylabel={FP / cow-day},
  symbolic x coords={auto,heatwave,snow,rain},
  xtick=data, legend style={font=\scriptsize, at={(0.5,1.25)}, anchor=north,
  legend columns=3}, ymin=0, enlarge x limits=0.15,
]
\addplot table[x=regime,y=telemetry] {data/expB_fp.dat};
\addplot table[x=regime,y=self] {data/expB_fp.dat};
\addplot table[x=regime,y=social] {data/expB_fp.dat};
\legend{telemetry,+self,+social}
\end{axis}
\end{tikzpicture}
\caption{False-positive rate by weather regime and feature family. Near-zero
throughout; the social features do not inflate false alarms.}
\label{fig:fp}
\end{figure}

\paragraph{Detection latency and ablation (C, D).} Table~\ref{tab:lat} gives
median detection latency (hours) and detection rate by condition and feature
family. All three conditions are detected in every replicate (100\% detection)
within the observation window. Contrary to our initial expectation---and to an
earlier single-run observation that motivated the social features---the ablation
returns a \emph{null} result: adding self-relative and social features does not
reduce median detection latency for any condition, and marginally increases it
for lameness (consistent with mild overfitting as feature count grows). Illness
is already identified at \illlatsocial\,h by both the telemetry-only and full
models. We report this openly; it is the kind of finding the single-run
evaluations this harness replaces would have hidden.

\begin{table}[t]\centering
\caption{Median detection latency (h) and detection rate (\%) by condition and
feature family, over \numreplicates{} replicates.}
\label{tab:lat}
\begin{tabular}{lrrrrrr}
\toprule
& \multicolumn{2}{c}{Telemetry} & \multicolumn{2}{c}{+Self} & \multicolumn{2}{c}{+Social} \\
\cmidrule(lr){2-3}\cmidrule(lr){4-5}\cmidrule(lr){6-7}
Condition & h & \% & h & \% & h & \% \\
\midrule
ill & 2.4 & 100 & 2.4 & 100 & 2.4 & 100 \\
lame & 4.0 & 100 & 4.0 & 100 & 4.5 & 100 \\
oestrus & 2.8 & 100 & 2.8 & 100 & 2.8 & 100 \\
\bottomrule
\end{tabular}

\end{table}

\section{Discussion}
The strong result is robustness. \emph{Herd-relative} normalisation delivers the
near-zero false-alarm rate in Table~\ref{tab:fp} even under sustained forced
weather: because every feature is a $z$-score against the concurrent herd, a
covariate that moves the herd together cancels. This is the property that
threshold-based systems lack and the one most likely to transfer to the field.

The ablation is a cautionary result and, we think, an instructive one. We
introduced self- and social-relative features on the \emph{change-not-state}
principle---a constitutionally slow or peripheral animal is not the same as one
that has \emph{become} so---and an early single run suggested they accelerated
illness detection. Under \numreplicates{} seeded replicates that effect
disappears (Table~\ref{tab:lat}). The mechanism is diagnosable: the illness
signature includes spatial withdrawal, which the herd-relative feature set
already captures through distance from the herd centroid, and the acute
fever/rumination collapse is fast and unambiguous---so the network-strength
features are largely redundant \emph{for detection latency in this simulator}.
Their remaining value here is interpretability (an explicit withdrawal signal in
the alert) and situational awareness (the live sociogram), not earlier alarms.

We would expect the social features to matter more where telemetry is noisier or
withdrawal is not already encoded spatially---for example under GNSS error and
duty-cycling, or for conditions whose primary early sign is social rather than
spatial. Testing that requires real data, and is precisely the kind of
hypothesis this benchmark is meant to generate rather than settle.

\section{Limitations}
\label{sec:limitations}
\textbf{No real-herd validation.} The single most important caveat: all results
are in-silico. The simulator's condition signatures are the same signals the
detector keys on, so absolute accuracies here upper-bound what noisy real data
would yield and must not be read as field performance. \textbf{Shared-assumption
risk.} Simulator and detector share an author and a worldview; independent data
is required to falsify the approach. \textbf{Proximity $\neq$ contact.} GNSS
proximity is a coarse proxy for social association; real deployments face fix
error and duty-cycling. \textbf{Behaviour-class noise.} We assume a perfect
activity classifier; real classifiers are $\sim$90--95\% accurate.
\textbf{Simple model and narrow label set.} A linear classifier over three
conditions is deliberately minimal; richer models and more conditions (calving,
metabolic disease, co-morbidities) remain untested.

\section{Future work}
The essential next step is a field study pairing collar telemetry with
veterinary diagnoses, on which the same pipeline can be validated and the
simulator calibrated or replaced. Methodologically: label noise on the activity
classifier; richer network features (betweenness, community stability); directed
approach patterns for oestrus; and sequence models replacing the linear
classifier once real data justifies the capacity.

\section{Conclusion}
Taking the herd as the unit of normalisation, and reading the social network
collars already sample, yields a detection method that is robust to herd-wide
covariates and sensitive to individual change---demonstrated here on a
reproducible simulation benchmark. We release the simulator, detector, and
harness to invite replication and, above all, field validation.

{\small
\bibliographystyle{plainnat}
\bibliography{references}
}

\end{document}
